\section{The Narrow Gate}

Our Lord was asked once how many would be saved, and his answer is recorded by Saint Luke. ``And a certain man said to him: Lord, are they few that are saved? But he said to them: Strive to enter by the narrow gate: for many, I say to you, shall seek to enter and shall not be able. But when the master of the house shall be gone in and shall shut the door, you shall begin to stand without; and knock at the door, saying: Lord, open to us. And he answering, shall say to you: I know you not, whence you are.''\endnote{Luke 13:23--25, Douay--Rheims (Challoner).} In the Sermon on the Mount he had said the same without being asked:

\begin{quote}
Enter ye in at the narrow gate: for wide is the gate, and broad is the way that leadeth to destruction, and many there are who go in thereat. How narrow is the gate, and strait is the way that leadeth to life: and few there are that find it!\endnote{Matthew 7:13--14, Douay--Rheims (Challoner).}
\end{quote}

\noindent And at the end of the parable of the wedding feast, after the guest without a garment has been cast into the exterior darkness, he closed with this sentence: ``For many are called, but few are chosen.''\endnote{Matthew 22:14, Douay--Rheims (Challoner).}

\subsection{How the Fathers and Doctors read it}

Chrysostom, Augustine, Gregory, Thomas and Alphonsus take these words as they stand: the way to life is narrow, and those who walk in it are fewer than those who do not. Chrysostom, preaching on the narrow gate, told his people that in saying, ``Few there be that find it,'' the Lord exposed the carelessness of most men and ``instructed His hearers not to regard the felicities of the many, but the labors of the few. For the more part, saith He, so far from walking this way, do not so much as make it their choice.''\endnote{Chrysostom, \emph{Homilies on Matthew} 23.7.} Later, preaching on the Acts of the Apostles, he warned his people not to trust in their numbers: ``Many more in number are they who go down into hell: but greater than it is the Kingdom, however few it contain.'' Then he asked how many in his city were likely to be saved, and answered his own question: ``It is disagreeable, what I am going to say, but I will say it nevertheless. Among all these myriads, there are not to be found one hundred likely to be saved: nay, even as to these, I question it.''\endnote{Chrysostom, \emph{Homilies on the Acts of the Apostles} 24, \emph{Nicene and Post-Nicene Fathers}, 1st series, vol.~11.}

Augustine, reflecting on the mass of mankind wounded by the first sin and on the mercy that draws some out of it, wrote that ``many more are left under punishment than are delivered from it, in order that it may thus be shown what was due to all.''\endnote{Augustine, \emph{City of God} XXI.12.} Gregory the Great, preaching on the labourers in the vineyard, stopped at the words \latin{Multi enim sunt vocati, pauci vero electi} and turned them on his own congregation: many come to faith, he said, and few are brought to the heavenly kingdom; and looking at the walls of the church filled for the feast, he asked who could know how few of all that crowd were numbered in the flock of God's elect, \latin{quis sciat quam pauci sunt qui in illo electorum Dei grege numerentur?}\endnote{Gregory the Great, \emph{Homilies on the Gospels} 19.5, Latin text, Wikisource, \emph{Homiliarum in Evangelia}, accessed 16 October 2026.} In another homily, on the wedding garment, he drew a lesson for every Christian, \latin{Quia vocati sumus, novimus; si sumus electi, nescimus}: each of us knows that he has been called, and none knows whether he is chosen.\endnote{Gregory, \emph{Homilies on the Gospels} 38.14.}

Saint Thomas, in the First Part of the \emph{Summa}, gives his reason. What exceeds the common state of nature is found in the few, as profound learning is found in the few; and eternal life with God exceeds all nature, the more so in a nature wounded by sin. Therefore ``those who are saved are in the minority. In this especially, however, appears the mercy of God, that He has chosen some for that salvation, from which very many in accordance with the common course and tendency of nature fall short.''\endnote{Aquinas, \emph{Summa theologiae} I, q.~23, a.~7, reply to objection 3, trans.\ Fathers of the English Dominican Province (Project Gutenberg eBook 17611).} And Alphonsus, warning those who sin and count on God's patience, states it as the common opinion of his own day: ``Were God to bear for ever with sinners, no one should be damned: but the most common opinion is, that the greater part of adults, even among Christians, are lost.''\endnote{Alphonsus, \emph{Preparation for Death}, Consideration XVII, first point, p.~208.}

\subsection{The number known to God}

The saints who spoke these hard words spoke them as a warning to each of their hearers. Chrysostom spoke so that his hearers would ``not regard the many, nor be troubled thereat, but emulate the few.''\endnote{Chrysostom, \emph{Homilies on Matthew} 23.7.} Gregory spoke so that each might fear for himself, not knowing whether he was chosen. And Aquinas, in the very article in which he says the saved are the minority, will not name a number. He records the guesses of others---that as many men will be saved as angels fell, or as many as remained---and sets them aside: ``It is, however, better to say that, `to God alone is known the number for whom is reserved eternal happiness.'\,''\endnote{Aquinas, \emph{Summa theologiae} I, q.~23, a.~7, corpus.} The words he quotes are the Church's own, from the Roman Missal's prayer for the living and the dead: \latin{Deus, cui soli cognitus est numerus electorum in superna felicitate locandus}, ``God, to Whom alone is known the number of the elect who shall attain the happiness of heaven.''\endnote{Roman Missal, Additional Prayers (\latin{Orationes diversae}), ``For the Living and the Dead,'' Secret, \latin{Deus, cui soli}; Latin and English in Lasance and Walsh, \emph{The New Roman Missal} (1945).}

The Church has defined no number. She prays that her children may pass through the narrow gate, and she directs the fear the Fathers taught toward the possibility of one's own loss rather than toward the multitude of the lost.

\subsection{The son of perdition}

One man the Lord singled out. At the Last Supper he said of his betrayer, ``woe to that man by whom the Son of man shall be betrayed. It were better for him, if that man had not been born.''\endnote{Matthew 26:24, Douay--Rheims (Challoner).} In his prayer to the Father that same night he said, ``Those whom thou gavest me have I kept: and none of them is lost, but the son of perdition: that the scripture may be fulfilled.''\endnote{John 17:12, Douay--Rheims (Challoner).} And the Church, at the Mass of Holy Thursday and again on Good Friday, prays in words that have stood in her Missal for centuries:

\begin{quote}
\latin{Deus, a quo et Judas reatus sui poenam, et confessionis suae latro praemium sumpsit, concede nobis tuae propitiationis effectum: ut, sicut in passione sua Jesus Christus, Dominus noster, diversa utrisque intulit stipendia meritorum; ita nobis, ablato vetustatis errore, resurrectionis suae gratiam largiatur.}

O God, from Whom Judas received the punishment of his crime and the thief the reward of his confession, grant us the effect of Thy clemency, that, as Jesus Christ, our Lord, in His passion dealt according to their deserts with the one and the other, so, putting away from us the error of the past, He may bestow upon us the grace of His resurrection.\endnote{1962 Missal, Thursday of the Lord's Supper, Collect, repeated on Good Friday after the first lesson. English from Lasance and Walsh, \emph{The New Roman Missal} (1945), Good Friday.}
\end{quote}

The collect names two men. Both were sinners, and both stood near Christ in his passion. One received the reward of his confession, ``Lord, remember me when thou shalt come into thy kingdom,'' and heard, ``This day thou shalt be with me in paradise.''\endnote{Luke 23:42--43, Douay--Rheims (Challoner).} The other received the punishment of his crime. Augustine saw where the difference lay. Judas too repented after his fashion: ``I have sinned in betraying innocent blood.''\endnote{Matthew 27:4, Douay--Rheims (Challoner).} But ``by despairing of God's mercy in his sorrow that wrought death, he left to himself no place for a healing penitence.''\endnote{Augustine, \emph{City of God} I.17.} Peter denied his Lord three times on the same night; then he remembered the word of Jesus, ``And going forth, he wept bitterly,''\endnote{Matthew 26:75, Douay--Rheims (Challoner).} and he was forgiven.

\subsection{Judgment belongs to God}

The Church raises her saints to the altars, declaring that they reign with Christ, but beyond what God has revealed she pronounces no sentence of damnation on any soul. ``Although we can judge that an act is in itself a grave offense,'' the \emph{Catechism} teaches, ``we must entrust judgment of persons to the justice and mercy of God.''\endnote{\emph{Catechism of the Catholic Church} 1861.} John Paul II said the same of hell itself: ``Damnation remains a real possibility, but it is not granted to us, without special divine revelation, to know which human beings are effectively involved in it.''\endnote{John Paul II, General Audience of 28 July 1999, n.~4, official Vatican English, vatican.va, accessed 16 October 2026.} The warning of the narrow gate is given so that each of us may strive to enter, and pray for all who are still on the way.
